\textbf{3.\ \textsc{Kirill on a swing} (8 points)} --- \textit{Kaarel Hänni.}
A swing has rigid rods attaching it to a horizontal pivot, so it can rotate
freely in a vertical plane. Initially the swing hangs vertically below the
pivot with Kirill standing upright on the seat; his friend gives him an initial
angular velocity $\omega_0$ about the pivot. Kirill is practising athletic
swinging: whenever the swing momentarily has zero angular velocity, he quickly
squats; whenever the rods are again vertical, he quickly stands up. Treat Kirill
as a point mass at distance $a$ from the pivot when standing and $b$ when
squatting ($a < b$); squatting and standing are motions along the rods. Neglect
friction, air resistance, and the mass of the swing. Gravitational acceleration
is $g$.

\textbf{i)} \textit{(1.5 points)} Sketch, qualitatively, how Kirill's angular
speed depends on time during the first full period of the swing.

\textbf{ii)} \textit{(2 points)} On the phase diagram (angular velocity vs.\
angle), sketch, qualitatively, the trajectory from the initial push until the
swing first goes over the top. The total number of periods shown need not be
correct.

\textbf{iii)} \textit{(4.5 points)} Find how many times Kirill must stand up
before the swing first goes over the top. Evaluate your answer for
$a = 2.5\,\mathrm{m}$, $b = 3.0\,\mathrm{m}$, $\omega_0 = 1.0\,\mathrm{rad/s}$,
and $g = 10\,\mathrm{m/s^2}$.
